\section{Additional Exercises}

\subsection{Penrose Diagram Of A Radiation-Filled Universe}

\begin{exercise}\label{ex:diagrams-1}
Find a differential equation for radial null geodesics in a spatially flat FRW universe filled with radiation, using the chart $(t,r,\vartheta,\phi)$ introduced in the lectures. Explicitly write down the precise range of the chart variables.
\end{exercise}
\hyperref[sol:diagrams-1]{\small \textit{Solution on p.~\pageref*{sol:diagrams-1}}}

\begin{exercise}\label{ex:diagrams-2}
Determine the $t$-coordinate of a geodesic in terms of the $r$ coordinate. Draw some of the null geodesics in the underlying chart.
\end{exercise}
\hyperref[sol:diagrams-2]{\small \textit{Solution on p.~\pageref*{sol:diagrams-2}}}

\begin{exercise}\label{ex:diagrams-3}
Find a chart in which the geodesics are lines of constant slope $\pm1$. Determine the range of the coordinates.
\end{exercise}
\hyperref[sol:diagrams-3]{\small \textit{Solution on p.~\pageref*{sol:diagrams-3}}}

\begin{exercise}\label{ex:diagrams-4}
Choose the so-called null coordinates $u$ and $v$ in which the null geodesics of positive slope are parallel to the $u$-axis and the ones of negative slope are parallel to the $v$-axis. Determine the range of the coordinates.
\end{exercise}
\hyperref[sol:diagrams-4]{\small \textit{Solution on p.~\pageref*{sol:diagrams-4}}}

\begin{exercise}\label{ex:diagrams-5}
Compactify, i.e., rescale to finite ranges, each of the two null coordinates by an appropriate transformation. Determine the range of the coordinates.
\end{exercise}
\hyperref[sol:diagrams-5]{\small \textit{Solution on p.~\pageref*{sol:diagrams-5}}}

\begin{exercise}\label{ex:diagrams-6}
By final transformation, recover the notion of temporal and radial coordinates. Determine the ranges of those coordinates. Draw the Penrose-Carter diagram.
\end{exercise}
\hyperref[sol:diagrams-6]{\small \textit{Solution on p.~\pageref*{sol:diagrams-6}}}
